\begin{statement}{C.2.UN.5}{Proved}{Every normal irreducible support is a smooth torus quotient}
Let $f:Y\to B$ be a proper holomorphic map from a smooth
complex threefold to a smooth curve, with connected fibres
and smooth general fibre a complex two-torus. Let $mS$ be
a full fibre whose reduced support is irreducible and normal.
Then $S$ is smooth and $e(S)=0$. More precisely, the
normalized base change $t=s^m$ gives a connected two-torus
$A$ and a finite unramified cyclic map
\[
                         A\longrightarrow S
\]
of degree $m$. Thus $S$ is a two-torus or a bielliptic
surface. In the bielliptic case the order $n$ of $K_S$ is
$2,3,4$, or $6$, and $n\mid m$; in particular a reduced
normal fibre is a torus.

Every two-torus occurs with every positive multiplicity.
Every bielliptic surface of canonical order $n$ occurs
with every multiplicity divisible by $n$. These
realizations can be made in smooth compact threefolds
over $\PP^1$.
\end{statement}
\begin{proof}
A normal surface hypersurface has only isolated
singularities. At a singular point of $S$, write the
actual map as $t=u h^m$, with $u$ a holomorphic unit
and $h$ a reduced local equation of $S$. A local
$m$th root of $u$ gives $t=(u^{1/m}h)^m$. The
isolated-hypersurface assertion of A.5.17 supplies an
embedded vanishing two-sphere of self-intersection
$-2$ in a small nonzero Milnor fibre of $u^{1/m}h$.
This fibre is an open part of the actual nearby
torus fibre of $f$. A sphere in a torus is nullhomotopic,
and hence has zero homology class and zero
self-intersection. The contradiction proves that $S$
is smooth. This argument uses no Betti-number or
K\"ahler hypothesis on the ambient threefold.

Shrink the base to a disc with coordinate $t$. Near
any point of $S$ write $t=u z^m$ with $z$ a smooth
equation of $S$. The normalization $U'$ of the base
change $t=s^m$ is given locally by the explicit maps
\[
        w^m=u,\qquad s=wz,
        \qquad p:U'\longrightarrow U.
 \tag{C.UNN.1}
\]
Since $u$ never vanishes, $p$ is finite and unramified
of degree $m$; these local normalizations glue to the
normalization of the actual fibre product. The map
$h:U'\to\Delta_s$ has $ds=w\,dz$ along its central
support, so that support is smooth and reduced.
For each $s\ne0$, the map $p$ identifies $h^{-1}(s)$
with $f^{-1}(s^m)$, a connected torus. The Stein
factor of $h$ has generic degree one over the normal
disc and is therefore the disc. In particular $U'$
and its central fibre $A$ are connected. After
shrinking the disc, $h$ is a proper submersion.
The deck action is $w\mapsto\zeta_m w$, so it is
free also on $A$, and its quotient there is exactly
$S$.

We verify that $A$ is a torus through its actual
Albanese map. Ehresmann's theorem identifies its
integral cohomology ring with that of a real
four-torus. Thus $b_1(A)=4$, and
\cite[Theorem~11]{Buchdahl1999} makes $A$ K\"ahler.
The period-integral map
$a:A\to\operatorname{Alb}(A)$ is consequently
defined by its two independent closed holomorphic
one-forms and induces an isomorphism on integral
$H^1$. Both surfaces have the exterior algebra on
this lattice as their integral cohomology ring.
The pullback of a top-degree generator is therefore
a generator, so $a$ has degree one; holomorphicity
fixes the positive sign. Its pullback on $H^2$ is
also an isomorphism. A curve contracted by $a$
would pair to zero with every pullback class and
hence with a K\"ahler class on $A$, contradicting
positive area. Thus $a$ is finite of degree one
and is an isomorphism. This proves the torus
identification, without imposing a metric on $U'$.
It also agrees with the deformation theorem
\cite[Theorem~4.6]{Catanese2002}.

Write a generator of the deck group on
$A=\mathbb C^2/\Lambda$ as $g(z)=Lz+b$.
Every holomorphic torus automorphism has this affine
form: the derivative of a lift is periodic and hence
constant on the compact quotient. The matrix $L$
preserves $\Lambda$ and has finite order. If $1$
were not an eigenvalue of $L$, the equation
$(I-L)z=b$ in $\mathbb C^2$ would give a fixed
point of $g$, contrary to the free deck action.
If $L=I$, the deck group acts by translations and
$S$ is the torus quotient by the resulting finite
subgroup.

Otherwise the eigenvalues on $\mathbb C^2$ are
$1$ and a nontrivial root of unity $\zeta$.
The corresponding integral real representation
has characteristic polynomial
$(x-1)^2(x-\zeta)(x-\bar\zeta)$, so the order
$n$ of $\zeta$ is $2,3,4$, or $6$. The images of
$I-L$ and of $I+L+\cdots+L^{n-1}$ are elliptic
subtori with complementary tangent lines. Their
addition map is an isogeny from a product of
elliptic curves to $A$. The kernel of the deck
group's linear representation consists of
translations. Quotienting by it first gives a
torus $A_1$ with a free cyclic action of order
$n$ and quotient $S$. On the canonical form
this action has character $\zeta$, so $K_S$
has exact order $n$, while the invariant
holomorphic one-form space has dimension one.
Equivalently, lifting to the product isogeny
and removing the pure translation kernels gives
the bielliptic presentation
$(E\times F)/G$, with $G$ acting faithfully
by translations on $E$ and with $F/G\simeq\PP^1$.
These are precisely the bielliptic surfaces
of \cite[Definition~3.1 and Theorems~3.2--3.3]{Catanese2002}.
The linear order divides the deck order $m$.
In both cases the unramified cover gives
$m\,e(S)=e(A)=0$.

We finish with realizations and their multiplicities.
Given any free cyclic torus quotient
$A\to S=A/\langle g\rangle$ of degree $m$, set
\[
 Y=(A\times\PP^1)/\langle(a,z)\mapsto
                         (g(a),\zeta_m z)\rangle,
 \qquad f([a,z])=z^m.
 \tag{C.UNN.2}
\]
The action is free, hence $Y$ is smooth and compact.
For $t\ne0,\infty$ its $m$ slices are permuted
transitively, and the quotient fibre is exactly
$A$. The full fibres at $0$ and $\infty$ are
$mS$. Thus these are actual compact realizations
of the displayed support and multiplicity.
For a prescribed torus $S$, an index-$m$ sublattice
with cyclic quotient gives the required cyclic
isogeny $A\to S$ for every $m\ge1$.

For a prescribed bielliptic surface, use its
product presentation $(E\times F)/G$. The linear
character $G\to\mu_n$ is surjective, and its
kernel acts by translations on both factors.
Consequently its canonical cyclic cover is a
torus. Moreover the character factors through
$\pi_1(E/G)\simeq\mathbb Z^2$, because $G$
acts on $E$ by translations. A surjective
character $\mathbb Z^2\to\mathbb Z/n$ lifts
to a primitive integral homomorphism
$\mathbb Z^2\to\mathbb Z$: choose integer
lifts of its two coefficients, and change one
by a multiple of $n$ to make their gcd one.
The possible prime divisors of the other
coefficient are avoided by the Chinese
remainder theorem; primes dividing $n$ already
cannot divide both coefficients. Pull this
homomorphism back along the connected elliptic
fibration $S\to E/G$ and reduce it modulo
$m=kn$. Its kernel defines a connected cyclic
unramified cover of degree $m$ that factors
through the canonical torus cover. A connected
finite unramified cover of a torus is the
quotient of $\mathbb C^2$ by a finite-index
sublattice, so this new cover is again a
torus. Applying (C.UNN.2) proves every claimed
bielliptic multiplicity.
\end{proof}


% C.2.UN.18 from part_C2_projective_elliptic.tex

